% ============================================================
%  RAPPORT — Introduction au Traitement de Données
%  « L'anatomie d'un grand café arabica »
%  MAM3 — Polytech Nice Sophia — 2025-2026
% ============================================================
\documentclass[12pt,a4paper]{article}

% ---- Packages ----
\usepackage[T1]{fontenc}
\usepackage[utf8]{inputenc}
\usepackage[french,english]{babel}
\usepackage{charter}
\usepackage[top=2.5cm,bottom=2.5cm,left=2.5cm,right=2.5cm]{geometry}
\usepackage{graphicx}
\usepackage{booktabs}
\usepackage{array}
\usepackage{xcolor}
\usepackage{colortbl}
\usepackage{amsmath,amssymb}
\usepackage{hyperref}
\usepackage{float}
\usepackage{subcaption}
\usepackage{enumitem}
\usepackage{titlesec}
\usepackage{fancyhdr}
\usepackage{microtype}
\usepackage{parskip}
\usepackage{mdframed}
\usepackage{tcolorbox}
\usepackage{wrapfig}
\usepackage{multirow}

% ---- Couleurs thème café ----
\definecolor{coffeebrown}{RGB}{74,44,42}
\definecolor{coffeemedium}{RGB}{123,79,46}
\definecolor{coffeecream}{RGB}{245,230,200}
\definecolor{coffeegold}{RGB}{212,168,67}
\definecolor{accentgreen}{RGB}{92,122,62}
\definecolor{lightgray}{RGB}{245,245,245}

% ---- Hyperref ----
\hypersetup{
  colorlinks=true,
  linkcolor=coffeebrown,
  citecolor=accentgreen,
  urlcolor=coffeemedium,
  pdftitle={L'anatomie d'un grand café arabica},
  pdfauthor={MAM3 — Polytech Nice Sophia}
}

% ---- En-têtes / Pieds de page ----
\pagestyle{fancy}
\fancyhf{}
\fancyhead[L]{\small\textcolor{coffeebrown}{\textit{Introduction au Traitement de Données — MAM3}}}
\fancyhead[R]{\small\textcolor{coffeebrown}{\thepage}}
\fancyfoot[C]{\small\textcolor{gray}{\textit{Polytech Nice Sophia — 2025-2026}}}
\renewcommand{\headrulewidth}{0.4pt}
\renewcommand{\headrule}{\hbox to\headwidth{\color{coffeebrown}\leaders\hrule height \headrulewidth\hfill}}

% ---- Titres de section ----
\titleformat{\section}
  {\Large\bfseries\color{coffeebrown}}
  {\thesection.}{0.5em}{}[\vspace{2pt}\color{coffeegold}\hrule height 1.5pt\vspace{4pt}]

\titleformat{\subsection}
  {\large\bfseries\color{coffeemedium}}
  {\thesubsection.}{0.5em}{}

\titleformat{\subsubsection}
  {\normalsize\bfseries\color{accentgreen}}
  {\thesubsubsection.}{0.5em}{}

% ---- Boîtes de mise en valeur ----
\tcbuselibrary{skins,breakable}
\newtcolorbox{insight}{
  enhanced,colback=coffeecream,colframe=coffeebrown,
  fonttitle=\bfseries\color{coffeebrown},
  title=Insight clé,arc=4pt,boxrule=1pt,
  left=6pt,right=6pt,top=4pt,bottom=4pt
}
\newtcolorbox{methode}{
  enhanced,colback=lightgray,colframe=accentgreen,
  fonttitle=\bfseries\color{accentgreen},
  title=Méthode,arc=4pt,boxrule=1pt,
  left=6pt,right=6pt,top=4pt,bottom=4pt
}

% ============================================================
\begin{document}
% ============================================================

% ---- Page de titre ----
\begin{titlepage}
\begin{center}
\vspace*{1cm}
{\color{coffeebrown}\rule{\linewidth}{2pt}}\\[0.4cm]
{\Huge\bfseries\color{coffeebrown} L'anatomie d'un grand café arabica}\\[0.3cm]
{\Large\color{coffeemedium}\textit{Quels critères définissent objectivement un café d'exception ?}}\\[0.4cm]
{\color{coffeebrown}\rule{\linewidth}{2pt}}\\[1.5cm]

\includegraphics[width=0.55\linewidth]{fig1_correlogram.png}\\[1cm]

\begin{tabular}{rl}
\textbf{\color{coffeebrown}Cours :} & Introduction au Traitement de Données\\
\textbf{\color{coffeebrown}Formation :} & MAM3 — Mathématiques Appliquées et Modélisation\\
\textbf{\color{coffeebrown}Établissement :} & Polytech Nice Sophia — Université Côte d'Azur\\
\textbf{\color{coffeebrown}Année :} & 2025-2026\\
\textbf{\color{coffeebrown}Dataset :} & CQI Arabica Coffee Quality Dataset (Kaggle)\\
\textbf{\color{coffeebrown}Dépôt Git :} & \href{https://github.com/}{[URL du dépôt Git]}
\end{tabular}\\[2cm]

\end{center}
\end{titlepage}

% ---- Table des matières ----
\tableofcontents
\newpage

% ============================================================
\section{Business Goal — Le café, un objet de science}
% ============================================================

\subsection{Contexte et motivation}

Le café est l'une des boissons les plus consommées au monde, avec plus de 2,5 milliards de tasses
bues chaque jour. Pourtant, la frontière entre un café ordinaire et un café d'\emph{exception}
reste souvent floue pour le consommateur. Pour les professionnels du secteur — torréfacteurs,
importateurs, baristas — cette question est au cœur de décisions économiques majeures.

Le \textbf{Coffee Quality Institute (CQI)} propose un protocole d'évaluation sensorielle
rigoureux, le \emph{cupping score}, qui décompose la qualité d'un café en huit critères objectifs,
chacun noté de 6 à 10. Le score total (\textit{Total Cup Points}) est la somme de ces sous-scores.
Un café atteignant 80 points ou plus est qualifié de \emph{Specialty Coffee} — la crème de la crème.

\begin{insight}
  \textbf{Question de recherche :} Quels sont les facteurs — sensoriels, géographiques,
  variétaux et agronomiques — qui permettent de prédire et d'expliquer la qualité d'un café arabica
  noté par le CQI ?
\end{insight}

\subsection{Utilité du dataset}

Le dataset \texttt{df\_arabica\_clean.csv} (source : Kaggle, CQI) regroupe \textbf{207 échantillons}
de cafés arabica certifiés, notés entre 2021 et 2023, provenant de \textbf{23 pays}.
Il constitue une source idéale pour répondre à notre question car :
\begin{itemize}[noitemsep]
  \item Les notes sensorielles sont \textbf{standardisées} par des juges certifiés Q-Grader ;
  \item Il couvre une grande diversité géographique, variétale et de méthodes de traitement ;
  \item Il comprend à la fois des variables continues, catégorielles et des données de défauts,
        permettant une \textbf{ingénierie de features} riche.
\end{itemize}

% ============================================================
\section{Gestion d'équipe}
% ============================================================

\subsection{Organisation du travail}

Le travail a été organisé en phases successives, calées sur le calendrier du cours.

\begin{center}
\begin{tabular}{@{}lll@{}}
\toprule
\rowcolor{coffeecream}
\textbf{Phase} & \textbf{Tâche principale} & \textbf{Date} \\
\midrule
1 & Choix du dataset, formulation de la question & 01/06/2026 \\
2 & Nettoyage, wrangling, exploration & 01/06/2026 \\
3 & Visualisation et storytelling & 02/06/2026 \\
4 & Feature engineering, régression & 02--03/06/2026 \\
5 & Rédaction du rapport, préparation des slides & 04/06/2026 \\
6 & Soutenance orale & 05/06/2026 \\
\bottomrule
\end{tabular}
\end{center}

\subsection{Répartition des rôles}

\begin{itemize}[noitemsep]
  \item \textbf{Membre 1} — Chargé(e) de l'exploration des données, de la visualisation et
        du storytelling ;
  \item \textbf{Membre 2} — Chargé(e) du feature engineering et de la régression ;
  \item \textbf{Membre 3} — Chargé(e) de la rédaction du rapport et de la préparation des slides.
\end{itemize}

% ============================================================
\section{Description des données}
% ============================================================

\subsection{Vue d'ensemble}

Le dataset contient \textbf{207 lignes} et \textbf{41 colonnes}, couvrant trois familles de variables :

\begin{center}
\begin{tabular}{@{}lll@{}}
\toprule
\rowcolor{coffeecream}
\textbf{Famille} & \textbf{Variables clés} & \textbf{Type} \\
\midrule
\textbf{Identité} & Pays, Région, Ferme, Variété, Propriétaire & Catégorielle / texte \\
\textbf{Agronomique} & Altitude, Méthode de traitement, Couleur, Humidité & Numérique / catég. \\
\textbf{Sensorielle} & Aroma, Flavor, Aftertaste, Acidity, Body, Balance, & Numérique [6--10] \\
& Uniformity, Clean Cup, Sweetness, Overall, Total Cup & \\
\textbf{Défauts} & Category One Defects, Category Two Defects, Quakers & Entier \\
\bottomrule
\end{tabular}
\end{center}

\subsection{Statistiques descriptives des critères sensoriels}

\begin{center}
\begin{tabular}{@{}lrrrrr@{}}
\toprule
\rowcolor{coffeecream}
\textbf{Critère} & \textbf{Min} & \textbf{Q1} & \textbf{Médiane} & \textbf{Moyenne} & \textbf{Max} \\
\midrule
Aroma        & 6.50 & 7.58 & 7.67 & 7.72 & 8.58 \\
Flavor       & 6.75 & 7.58 & 7.75 & 7.74 & 8.50 \\
Aftertaste   & 6.67 & 7.42 & 7.58 & 7.60 & 8.42 \\
Acidity      & 6.83 & 7.50 & 7.67 & 7.69 & 8.58 \\
Body         & 6.83 & 7.50 & 7.67 & 7.64 & 8.25 \\
Balance      & 6.67 & 7.50 & 7.67 & 7.64 & 8.42 \\
\midrule
\textbf{Total Cup Points} & 78.00 & 82.58 & 83.75 & 83.71 & 89.33 \\
\bottomrule
\end{tabular}
\end{center}

\paragraph{Note :} Les critères \textit{Uniformity}, \textit{Clean Cup} et \textit{Sweetness}
présentent des variances quasi nulles (presque tous égaux à 10 dans ce dataset premium), ce qui
explique leur faible pouvoir prédictif.

\subsection{Nettoyage et wrangling}

\begin{itemize}[noitemsep]
  \item L'\textbf{altitude} est fournie comme texte libre ; elle a été convertie en numérique
        (\texttt{pd.to\_numeric(..., errors=`coerce')}), aboutissant à 153 valeurs valides sur 207.
  \item Les \textbf{méthodes de traitement} rares ($n = 1$) ont été regroupées en une catégorie
        « Autres » pour les analyses comparatives.
  \item La \textbf{couleur} des grains présente des incohérences typographiques (\textit{yellow green}
        vs \textit{yellow-green}) ; un nettoyage par normalisation de chaînes a été appliqué.
  \item \textbf{Aucune ligne n'a été supprimée} pour les analyses sensorielles (207 valeurs complètes).
\end{itemize}

% ============================================================
\section{Visualisation des données}
% ============================================================

\subsection{Corrélations entre les critères sensoriels}

\begin{figure}[H]
  \centering
  \includegraphics[width=0.72\textwidth]{fig1_correlogram.png}
  \caption{Corrélogramme triangulaire inférieur des huit critères sensoriels (n\,=\,207).
           Les couleurs chaudes indiquent une forte corrélation positive.}
  \label{fig:corr}
\end{figure}

La figure~\ref{fig:corr} révèle que tous les critères sensoriels sont \textbf{fortement corrélés
entre eux} ($r > 0.70$ pour la plupart), et en particulier avec le score total.
Les corrélations les plus élevées avec le \textit{Total Cup Points} sont :
\textit{Overall} ($r = 0.95$), \textit{Flavor} ($r = 0.94$), \textit{Aftertaste} ($r = 0.93$)
et \textit{Balance} ($r = 0.93$).

\begin{insight}
  Un café arabica de qualité est un équilibre : aucun critère isolé ne suffit. La force du
  \textit{Flavor} et de l'\textit{Aftertaste} sont cependant les signaux les plus discriminants.
\end{insight}

\subsection{Profil sensoriel selon la méthode de traitement}

\begin{figure}[H]
  \centering
  \includegraphics[width=0.62\textwidth]{fig2_radar.png}
  \caption{Diagramme radar des profils sensoriels moyens selon la méthode de traitement.
           Les cafés « Natural/Dry » se distinguent par un Aroma et un Body plus prononcés.}
  \label{fig:radar}
\end{figure}

La figure~\ref{fig:radar} compare trois méthodes : voie humide (\textit{Washed}),
voie sèche (\textit{Natural}) et processus intermédiaire (\textit{Honey/Pulped}).
On observe que les cafés traités en voie sèche tendent à présenter un \textbf{arôme et un corps
plus intenses}, tandis que les cafés lavés sont plus équilibrés sur l'ensemble des critères.

\subsection{Scores par pays d'origine}

\begin{figure}[H]
  \centering
  \includegraphics[width=0.90\textwidth]{fig3_countries.png}
  \caption{Distribution du \textit{Total Cup Points} pour les 12 pays les plus représentés.
           Les boîtes à moustaches sont ordonnées par médiane croissante. La ligne pointillée
           indique la moyenne globale (83,7 pts).}
  \label{fig:pays}
\end{figure}

L'\textbf{Éthiopie} et la \textbf{Tanzanie} dominent le classement en termes de score médian
(respectivement 85,0 et 84,7 pts), loin devant le Brésil et El Salvador qui peinent à dépasser
82 pts en médiane. Taiwan, très représenté ($n=61$), affiche une médiane élevée (84,4 pts) avec une
grande dispersion, témoignant d'une filière dynamique mais hétérogène.

\subsection{Influence de l'altitude}

\begin{figure}[H]
  \centering
  \includegraphics[width=0.90\textwidth]{fig4_altitude.png}
  \caption{\textit{Gauche :} nuage de points altitude/score avec droite de régression.
           \textit{Droite :} score moyen (± IC à 95\,\%) par tranche d'altitude.
           La tendance positive est faible mais régulière au-delà de 1600\,m.}
  \label{fig:alt}
\end{figure}

La corrélation brute altitude-score est modeste ($r = 0.10$), mais la figure~\ref{fig:alt}
révèle un effet de seuil : les cafés cultivés \textbf{au-delà de 1800\,m} obtiennent en moyenne
\textbf{84,8 pts}, soit +1,3 pt par rapport aux cafés cultivés sous 1000\,m (83,5 pts).
La faible corrélation s'explique par la multiplicité des facteurs confondants (pays, variété,
méthode de traitement).

% ============================================================
\section{Feature Engineering}
% ============================================================

\subsection{Variables directement disponibles (features brutes)}

Les six critères sensoriels (\textit{Aroma, Flavor, Aftertaste, Acidity, Body, Balance}) constituent
nos features principales. Ils sont tous sur la même échelle [6--10] et normalement distribués,
ce qui facilite l'application d'une régression linéaire.

\subsection{Features construites}

\subsubsection{Score de qualité sensorielle globale}

Inspiré du pipeline de feature engineering du cours (Lecture 4), nous construisons un
\textbf{indice de qualité composite} :

\begin{equation}
  \text{Sensory\_Score} = \frac{\text{Flavor} + \text{Aftertaste} + \text{Balance}}{3}
\end{equation}

Ce triplet concentre les critères les plus corrélés au score final et permet une réduction
de dimensionnalité sans perte d'information majeure.

\subsubsection{Pénalité de défauts}

En nous inspirant de l'exemple de détection de fraude bancaire (Lecture 4 : $z = x_1^2 + x_2^2$),
nous créons une variable synthétique combinant les deux types de défauts :

\begin{equation}
  \text{Defect\_Penalty} = \text{Cat1\_Defects} \times 4 + \text{Cat2\_Defects} + \text{Quakers}
\end{equation}

La pondération (×4 pour les défauts de catégorie 1) reflète leur impact sensoriel plus sévère.

\subsubsection{Encodage de l'altitude}

L'altitude brute est convertie en une variable ordinale \textit{Altitude\_Band} (6 niveaux),
puis one-hot encodée, conformément aux recommandations du cours (Lecture 2) sur les données
ordinales à plusieurs catégories.

\subsubsection{Encodage de la méthode de traitement}

Les trois méthodes principales (\textit{Washed}, \textit{Natural}, \textit{Honey/Pulped})
sont encodées par one-hot encoding (Lecture 2), les méthodes rares ($n=1$) étant regroupées
dans une catégorie « Autres ».

\subsection{Vérification de la pertinence des features}

\begin{figure}[H]
  \centering
  \begin{subfigure}[b]{0.48\textwidth}
    \includegraphics[width=\textwidth]{fig5_varieties.png}
    \caption{Score par variété (violon)}
    \label{fig:var}
  \end{subfigure}
  \hfill
  \begin{subfigure}[b]{0.48\textwidth}
    \includegraphics[width=\textwidth]{fig7_defects.png}
    \caption{Impact des défauts sur le score}
    \label{fig:def}
  \end{subfigure}
  \caption{Pertinence de deux features construites. La variété \textit{Gesha} se démarque
           clairement (médiane $\approx$ 85,4 pts). Les défauts de catégorie 2 et les
           \textit{Quakers} montrent une corrélation négative marquée ($r = -0.32$ et $r = -0.32$).}
\end{figure}

% ============================================================
\section{Régression linéaire et logistique}
% ============================================================

\subsection{Régression linéaire : prédire le score total}

\subsubsection{Modèle}

Nous appliquons une régression linéaire multiple (Lecture 4) :

\begin{equation}
  \widehat{\text{TCP}} = \beta_0 + \beta_1 \cdot \text{Aroma} + \beta_2 \cdot \text{Flavor}
  + \beta_3 \cdot \text{Aftertaste} + \beta_4 \cdot \text{Acidity}
  + \beta_5 \cdot \text{Body} + \beta_6 \cdot \text{Balance} + \varepsilon
\end{equation}

Les variables sont standardisées (\textit{StandardScaler}) avant l'estimation pour pouvoir
comparer les coefficients.

\subsubsection{Résultats}

\begin{center}
\begin{tabular}{@{}lrr@{}}
\toprule
\rowcolor{coffeecream}
\textbf{Critère} & \textbf{Coefficient standardisé} & \textbf{Interprétation} \\
\midrule
Flavor      & 0.345 & Critère le plus influent \\
Aroma       & 0.333 & \\
Aftertaste  & 0.331 & \\
Balance     & 0.322 & \\
Acidity     & 0.309 & \\
Body        & 0.252 & Critère le moins influent \\
\midrule
Intercept   & 83.71 & Score moyen de référence \\
\bottomrule
\end{tabular}
\end{center}

\begin{figure}[H]
  \centering
  \includegraphics[width=0.72\textwidth]{fig6_regression.png}
  \caption{Importance relative des critères sensoriels (coefficients standardisés de la régression
           linéaire). Un coefficient plus élevé signifie qu'une augmentation d'un écart-type du
           critère entraîne une plus forte hausse du score total.}
  \label{fig:reg}
\end{figure}

\textbf{Performance du modèle :}
\begin{itemize}[noitemsep]
  \item $R^2 = 0.993$ (sur données d'entraînement) — le modèle explique 99,3\,\% de la variance ;
  \item $R^2_{\text{CV}} = 0.676$ (validation croisée 5-fold) — performance robuste hors-échantillon ;
  \item $\text{RMSE} = 0.149$ point — l'erreur de prédiction est inférieure à 0,15 point.
\end{itemize}

\begin{methode}
  La différence entre $R^2 = 0.99$ (in-sample) et $R^2_{\text{CV}} = 0.68$ (out-of-sample)
  s'explique par le fait que le \textit{Total Cup Points} est quasi-mécaniquement la somme des
  sous-scores dans le protocole CQI. Le modèle de régression est donc une quasi-tautologie :
  il capture la structure même du scoring CQI. Pour prédire la qualité d'un café \textit{avant
  le cupping} (à partir de l'altitude, du pays, de la variété), un $R^2$ de 0.68 est en revanche
  réaliste.
\end{methode}

\subsection{Régression logistique : identifier les cafés d'exception}

\subsubsection{Définition de la cible}

Nous définissons un café comme \textbf{d'exception} si son score dépasse 84 points :

\begin{equation}
  Y = \mathbf{1}[\text{Total Cup Points} \geq 84]
\end{equation}

Sur le dataset, 93 cafés (44,9\,\%) atteignent ce seuil.

\subsubsection{Modèle LOGIT}

Conformément au cours (Lecture 4) :

\begin{equation}
  \ln\left(\frac{P(Y=1|X)}{1-P(Y=1|X)}\right)
  = \beta_0 + \beta_1 \cdot \text{Aroma} + \cdots + \beta_6 \cdot \text{Balance}
\end{equation}

\subsubsection{Résultats}

\begin{center}
\begin{tabular}{@{}lr@{}}
\toprule
\rowcolor{coffeecream}
\textbf{Critère} & \textbf{Coefficient logistique} \\
\midrule
Aftertaste  & 1.804 \\
Acidity     & 1.615 \\
Aroma       & 1.546 \\
Body        & 1.399 \\
Flavor      & 1.189 \\
Balance     & 1.186 \\
\bottomrule
\end{tabular}
\end{center}

\textbf{Performance :} Précision en validation croisée 5-fold = \textbf{89.9\,\%}.
L'\textit{Aftertaste} et l'\textit{Acidity} sont les deux critères qui permettent le mieux de
distinguer un café d'exception d'un café ordinaire.

% ============================================================
\section{Conclusion}
% ============================================================

\subsection{Enseignements de l'analyse}

Notre analyse du dataset CQI Arabica apporte des réponses concrètes à la question
\textit{« Quels critères définissent objectivement un café d'exception ? »} :

\begin{enumerate}
  \item \textbf{Le Flavor et l'Aftertaste} sont les critères sensoriels les plus déterminants :
        c'est dans la persistance et la complexité des arômes que se joue l'excellence.
  \item \textbf{L'origine géographique} compte : l'Éthiopie et la Tanzanie produisent
        structurellement des cafés plus notés, en lien avec leurs terroirs (altitude, variétés
        locales, pratiques de fermentation).
  \item \textbf{La variété} \textit{Gesha} se démarque avec un score moyen de 85,4 pts,
        justifiant son statut de café premium sur les marchés internationaux.
  \item \textbf{L'altitude élevée} ($> 1800$\,m) est associée à des scores supérieurs,
        mais l'effet reste modéré et masqué par d'autres facteurs.
  \item \textbf{Les défauts} (Quakers, défauts de catégorie 2) ont un impact négatif mesurable
        ($r \approx -0.32$) et constituent une variable prédictive utile.

\end{enumerate}

\subsection{Limites et perspectives}

\begin{itemize}[noitemsep]
  \item Le dataset ne compte que 207 observations, ce qui limite la puissance statistique,
        notamment pour les analyses par pays ou variété.
  \item Les notes CQI sont une mesure subjective (jury de Q-Graders) ; leur variabilité
        inter-juges n'est pas capturée.
  \item Une extension naturelle serait d'ajouter des données de \textbf{prix de marché}
        pour tester si le score CQI prédit effectivement la valeur commerciale du café.
  \item Un modèle plus sophistiqué (Random Forest, XGBoost) pourrait améliorer la prédiction
        à partir des seules variables agronomiques (pays, altitude, variété).
\end{itemize}

% ============================================================
\section*{Références}
% ============================================================
\addcontentsline{toc}{section}{Références}

\begin{enumerate}
  \item Fillatre, L. (2025-2026). \textit{Introduction to Data Processing — Lectures 1--4}.
        Polytech Nice Sophia, Université Côte d'Azur.
  \item Coffee Quality Institute (CQI). \textit{Arabica Coffee Quality Dataset}.
        Kaggle. \url{https://www.kaggle.com/datasets}
  \item Wilke, C.O. (2019). \textit{Fundamentals of Data Visualisation}.
        O'Reilly Media. \url{https://clauswilke.com/dataviz/}
  \item Pedregosa, F. et al. (2011). Scikit-learn: Machine Learning in Python.
        \textit{Journal of Machine Learning Research}, 12, 2825--2830.
  \item Wexler, S., Shaffer, J., \& Cotgreave, A. (2017).
        \textit{The Big Book of Dashboards}. Wiley.
\end{enumerate}

\end{document}
